\begin{pastexam}[label={px:ncku-106-ee-ds-05}]{成大電機所 106 學年度 簡答題第 1 題}
（簡答）

\textbf{(20 pts)} Give a sequence of numbers 6, 2, 4, 7, 3, 9, 5 recorded by an array (the first index of the
array is zero).
\begin{enumerate}[label=\alph*.]
\item \textbf{(6 pts)} Sort the numbers in the increasing order by the quick sort algorithm and show the status
of the array each time after you change the location of pivot (\textbf{take the last digit as pivot}). Each
time the array [0\ldots n-1] is divided into two subarrays whose indices are
[0 \ldots\ $\lfloor\frac{n-1}{2}\rfloor$] and [$\lfloor\frac{n-1}{2}\rfloor+1$ \ldots\ n-1], respectively,
and the subarray with small index is handled first.
\item \textbf{(4 pts)} Show the condition when the algorithm has to take the longest time according to the
above numbers. What is the time complexity if there exist n numbers and why?
\item \textbf{(5 pts)} Show the status of the heap in an array before you want to sort the number
\textbf{in the decreasing order.}
\item \textbf{(5 pts)} Based on the result of b, you delete the root and exchange the root with the number in
the array with the largest index. Please show the status of the array after you legal the new heap.
\end{enumerate}
\end{pastexam}

\begin{pastsoln}
a、b 小題屬於快速排序，c、d 小題屬於堆積排序，都在資結書第~14~章。

答案：a.\ 依序為 $[2, 4, 3, 5, 6, 9, 7]$、$[2, 3, 4, 5, 6, 9, 7]$、$[2, 3, 4, 5, 6, 7, 9]$；
b.\ 每次的基準都是子陣列的最大或最小值（例如輸入已經排好），$\Theta(n^2)$；
c.\ 最小堆積 $[2, 3, 4, 7, 6, 9, 5]$；d.\ $[3, 5, 4, 7, 6, 9, 2]$。

\textbf{題意說明。} a 小題那句「每次把 $[0..n-1]$ 分成 $[0..\lfloor (n-1)/2 \rfloor]$ 和
$[\lfloor (n-1)/2 \rfloor + 1..n-1]$」照字面是從正中間切，那是合併排序的切法；快速排序是在基準最後的位置切開。
這組資料用「取最後一個當基準」的分割，三次分割的基準剛好都落在子陣列的正中間位置
$\lfloor (n-1)/2 \rfloor$，所以兩種說法得到同樣的切法。 下面照快速排序做：基準放到最後的位置之後就不再
移動，再遞迴處理它左邊和右邊，左邊（索引小的）先做。 d 小題的「result of b」應該是「result of c」（b 小題
沒有堆積），這裡照 c 的結果做。

\textbf{a.} 用 Lomuto 分割（資結書虛擬碼~14.8），以子陣列最後一個數 $x$ 為基準：$i$ 指著「$\le x$ 區」的
最後一格，$j$ 由左往右掃，$A[j] \le x$ 就把它換進 $\le x$ 區；掃完再把基準換到 $\le x$ 區的後面。
\begin{itemize}
\item \textbf{第 1 次分割}，$A[0..6] = [6, 2, 4, 7, 3, 9, 5]$，基準 $5$。 $j = 0$：$6 > 5$；$j = 1$：$2 \le 5$，和
$A[0]$ 交換得 $[2, 6, 4, 7, 3, 9, 5]$；$j = 2$：$4 \le 5$，和 $A[1]$ 交換得 $[2, 4, 6, 7, 3, 9, 5]$；$j = 3$：$7 > 5$；
$j = 4$：$3 \le 5$，和 $A[2]$ 交換得 $[2, 4, 3, 7, 6, 9, 5]$；$j = 5$：$9 > 5$。 最後基準和 $A[3]$ 交換：
\[
[2, 4, 3, \underline{5}, 6, 9, 7],
\]
基準 $5$ 在索引 $3 = \lfloor 6/2 \rfloor$。 接著處理 $A[0..2]$，再處理 $A[4..6]$。
\item \textbf{第 2 次分割}，$A[0..2] = [2, 4, 3]$，基準 $3$。 $j = 0$：$2 \le 3$，和自己交換；$j = 1$：$4 > 3$。
基準和 $A[1]$ 交換：
\[
[2, \underline{3}, 4, 5, 6, 9, 7],
\]
基準在索引 $1$。 兩邊 $[2]$、$[4]$ 都只有一個數，不用再做。
\item \textbf{第 3 次分割}，$A[4..6] = [6, 9, 7]$，基準 $7$。 $j = 4$：$6 \le 7$，和自己交換；$j = 5$：$9 > 7$。
基準和 $A[5]$ 交換：
\[
[2, 3, 4, 5, 6, \underline{7}, 9],
\]
基準在索引 $5$。 兩邊都只剩一個數，排序完成。
\end{itemize}
三次分割共比較 $6 + 2 + 2 = 10$ 次。

\textbf{b.} 快速排序最慢的情況是\textbf{每一次分割的基準都是子陣列的最大值或最小值}，分割後一邊是空的、
另一邊有 $m - 1$ 個數。 以這七個數為例，輸入已經由小到大排好 $[2, 3, 4, 5, 6, 7, 9]$（每次基準都是最大值）
或由大到小排好 $[9, 7, 6, 5, 4, 3, 2]$（基準輪流是最小、最大）時，比較次數都是
$6 + 5 + 4 + 3 + 2 + 1 = 21$ 次，是七個數所有排列裡最多的；只要每一層的最後一個數都是最大或最小，例如
$[6, 4, 3, 5, 2, 7, 9]$，也是 $21$ 次。

$n$ 個數時，分割長度 $m$ 的子陣列要比較 $m - 1$ 次。 最壞情況的遞迴式
\[
T(n) = T(n - 1) + T(0) + \Theta(n), \qquad
\text{比較次數} = (n - 1) + (n - 2) + \cdots + 1 = \frac{n(n - 1)}{2},
\]
所以是 $\Theta(n^2)$（資結書第~14.5~節）。 原因是每一層只把子陣列縮小一個，遞迴深度是 $n - 1$，而不是
平衡切法的 $\lg n$。

\textbf{c.} 要用堆積排序排出\textbf{由大到小}，用\textbf{最小堆積}：每次把根（目前最小的數）換到陣列最後面，
最後陣列就是由大到小。 （由小到大用最大堆積，是資結書虛擬碼~14.13 的寫法，方向相反而已。）

陣列從索引 $0$ 開始，$A[i]$ 的子節點是 $A[2i + 1]$、$A[2i + 2]$。 由下而上建堆積（資結書第~8.2~節），從最後一個
有子節點的 $i = \lfloor 7/2 \rfloor - 1 = 2$ 往回做 sift-down：
\begin{itemize}
\item $i = 2$：$A[2] = 4$，子節點 $9$、$5$，$4$ 最小，不動。
\item $i = 1$：$A[1] = 2$，子節點 $7$、$3$，$2$ 最小，不動。
\item $i = 0$：$A[0] = 6$，子節點 $2$、$4$，和較小的 $2$ 交換，得 $[2, 6, 4, 7, 3, 9, 5]$；$6$ 到了索引 $1$，子節點
$7$、$3$，和 $3$ 交換，得 $[2, 3, 4, 7, 6, 9, 5]$；$6$ 到了索引 $4$，沒有子節點，停。
\end{itemize}
排序前的最小堆積是
\[
[2, 3, 4, 7, 6, 9, 5].
\]
（逐一插入建堆積也得到同一個陣列：依序插入 $6, 2, 4, 7, 3, 9, 5$，只有插入 $2$ 時和 $6$ 交換、插入 $3$ 時和 $6$
交換。）

\textbf{d.} 把根 $2$ 和最後一格 $A[6] = 5$ 交換，$2$ 留在 $A[6]$ 不再屬於堆積：
\[
[5, 3, 4, 7, 6, 9 \mid 2].
\]
在 $A[0..5]$ 對根做 sift-down：$5$ 的子節點 $3$、$4$，和 $3$ 交換得 $[3, 5, 4, 7, 6, 9 \mid 2]$；$5$ 在索引 $1$，
子節點 $7$、$6$ 都比 $5$ 大，停。 結果
\[
[3, 5, 4, 7, 6, 9, 2],
\]
前六格是合法的最小堆積，最後一格是已經就位的 $2$。

如果 c 小題用最大堆積，排出來會是由小到大，和題目要的方向相反；最大堆積是 $[9, 7, 6, 2, 3, 4, 5]$，這不是
本題的答案。
\end{pastsoln}
